\documentclass[../DoAn.tex]{subfiles}
\begin{document}

\label{chap:phuong_phap_de_xuat}

Chương này trình bày phương pháp xây dựng hệ thống phát hiện drone dựa trên tín hiệu RF. Dựa trên cơ sở lý thuyết ở Chương~2, phương pháp được tổ chức thành một pipeline thống nhất: nhận dữ liệu I/Q, phân đoạn, tạo spectrogram, chuẩn hóa ảnh và phân loại bằng mô hình học sâu. Mục tiêu của chương là mô tả cách triển khai phương pháp; cấu hình thực nghiệm cụ thể và kết quả đánh giá được trình bày ở chương sau.

\section{Phát biểu bài toán và hướng tiếp cận}

Bài toán được phát biểu dưới dạng phân loại nhị phân trên từng đoạn tín hiệu RF. Đầu vào ban đầu là chuỗi mẫu I/Q thu từ nguồn dữ liệu công khai hoặc từ thiết bị SDR. Mỗi đoạn I/Q được chuyển thành ảnh spectrogram $\mathbf{S}$, sau đó đưa vào mô hình phân loại để dự đoán một trong hai nhãn:
\begin{equation}
    f_{\theta}: \mathbf{S} \longrightarrow \hat{y},
    \qquad
    \hat{y}\in\{0,1\},
\end{equation}
trong đó $f_{\theta}$ là mô hình có tham số $\theta$, nhãn $0$ tương ứng với \textit{Non-drone} và nhãn $1$ tương ứng với \textit{Drone}. Lớp \textit{Drone} biểu thị đoạn tín hiệu có chứa hoạt động RF của drone theo quy ước gán nhãn của dữ liệu. Lớp \textit{Non-drone} biểu thị đoạn không chứa hoạt động drone, có thể bao gồm nhiễu nền hoặc tín hiệu từ các nguồn phát khác. Cách phát biểu nhị phân phù hợp với mục tiêu chính của đồ án là phát hiện sự hiện diện của drone, thay vì nhận dạng chính xác chủng loại hoặc giao thức truyền thông.

Hướng tiếp cận của đồ án là không giải điều chế nội dung gói tin RF. Thay vào đó, hệ thống học các đặc trưng thời gian-tần số trên spectrogram, chẳng hạn burst, vùng băng thông chiếm dụng, vệt nhảy tần hoặc sự khác biệt giữa tín hiệu drone và nguồn phát nền. Để đánh giá mức độ phù hợp của mô hình tiền huấn luyện với miền spectrogram RF, đồ án khảo sát hai chiến lược: \textit{Feature Extraction-based Classification} và \textit{fine-tuning}. Đơn vị dự đoán chính trong phương pháp là một đoạn tín hiệu sau phân đoạn, mỗi spectrogram tương ứng với 50ms của bản ghi RF và được gán một nhãn duy nhất. Khi cần suy luận ở mức bản ghi hoặc phiên thu, các dự đoán theo đoạn có thể được tổng hợp ở bước sau. Cách tổ chức này giúp chuẩn hóa kích thước đầu vào, tăng số lượng mẫu huấn luyện và cho phép quan sát sự thay đổi quyết định theo thời gian.

\section{Quy trình tổng thể của phương pháp đề xuất}

Pipeline xử lý được thiết kế thống nhất cho cả dữ liệu huấn luyện, dữ liệu kiểm thử và dữ liệu RF tự thu. Luồng xử lý tổng quát gồm các bước:

\begin{figure}[htbp]
    \centering
    \begin{tikzpicture}[
        node distance=0.7cm and 0.55cm,
        box/.style={draw, rounded corners, minimum height=1cm, text width=2.35cm, align=center, fill=gray!8},
        arr/.style={-Latex, thick}
    ]
        \node[box] (rf) {Tín hiệu RF\\thu bằng SDR};
        \node[box, right=of rf] (iq) {Dữ liệu\\I/Q};
        \node[box, right=of iq] (prep) {Tiền xử lý\\và phân đoạn};
        \node[box, right=of prep] (spec) {Ảnh\\spectrogram};
        \node[box, right=of spec] (model) {Mô hình\\học sâu};
        \node[box, below=0.65cm of model] (out) {Drone /\\Non-drone};
        \draw[arr] (rf) -- (iq);
        \draw[arr] (iq) -- (prep);
        \draw[arr] (prep) -- (spec);
        \draw[arr] (spec) -- (model);
        \draw[arr] (model) -- (out);
    \end{tikzpicture}
    \caption{Luồng tổng quát của bài toán phát hiện drone dựa trên tín hiệu RF trong đồ án.}
    \label{fig:overall-problem}
\end{figure}
\FloatBarrier


Việc giữ cùng pipeline có vai trò quan trọng vì mô hình học trực tiếp từ hình học và phân bố cường độ của spectrogram. Nếu cấu hình xử lý ở giai đoạn suy luận khác với giai đoạn huấn luyện, đặc trưng đầu vào có thể thay đổi dù tín hiệu RF gốc tương tự nhau. Vì vậy, các tham số như kiểu dữ liệu I/Q, tốc độ lấy mẫu, độ dài đoạn, cấu hình STFT, quy tắc log scale, resize và chuẩn hóa ảnh cần được lưu cùng cấu hình mô hình.

Trong giai đoạn xây dựng mô hình, tập dữ liệu nhị phân được tạo bằng cách kết hợp mẫu \textit{Drone} từ DroneDetectv2 với mẫu \textit{Non-drone} do đồ án tự thu bổ sung. Hai nguồn dữ liệu này được xử lý bằng cùng pipeline tạo spectrogram và được chia thành tập huấn luyện, tập xác thực, tập kiểm thử. Tập huấn luyện dùng để tối ưu SVM hoặc \textit{fine-tuning} backbone; tập xác thực dùng để chọn siêu tham số và checkpoint; tập kiểm thử chỉ dùng để báo cáo kết quả sau cùng. Các tập RF drone tự thu dùng riêng cho đánh giá ngoài miền, không được dùng để cập nhật tham số.

Mỗi bước trong pipeline đều tạo ra dữ liệu trung gian có thể kiểm tra lại. Bản ghi I/Q giữ vai trò dữ liệu gốc; đoạn I/Q là đơn vị xử lý; spectrogram là đầu vào của mô hình; dự đoán theo đoạn là đầu ra để tính chỉ số. Việc lưu liên kết giữa các mức này giúp truy vết khi mô hình dự đoán sai, chẳng hạn từ một ảnh spectrogram quay lại bản ghi RF, vị trí thời gian, cấu hình STFT và chỉ báo chất lượng tín hiệu tương ứng.

\section{Thiết kế mô hình thu tín hiệu RF}

Phần thu tín hiệu được thiết kế như khối đầu vào của toàn bộ hệ thống phát hiện drone. Mục tiêu của khối này là thu nhận tín hiệu RF phát ra từ drone hoặc bộ điều khiển trong băng tần quan sát, chuyển tín hiệu cao tần về dạng tín hiệu băng gốc I/Q và lưu lại dưới dạng dữ liệu số phục vụ cho các bước xử lý tiếp theo. Trong nghiên cứu này, thiết bị SDR (Software-Defined Radio – vô tuyến định nghĩa bằng phần mềm) được sử dụng là \textit{bladeRF}. Thiết bị được kết nối với máy tính thông qua giao tiếp USB 3.0 và điều khiển bằng phần mềm SoapySDR để cấu hình tần số trung tâm, tốc độ lấy mẫu, băng thông và hệ số khuếch đại thu. Khối thu chỉ thực hiện số hóa tín hiệu RF và ghi lại dữ liệu I/Q, không thực hiện giải điều chế hay giải mã giao thức truyền thông của drone.

\begin{figure}[htbp]
    \centering
    \includegraphics[width=0.9\textwidth]{Hinh_ve/chap3/thu_tin_hieu.jpg}
    \caption{Bối cảnh thu tín hiệu trong thực nghiệm.}
    \label{fig:spectrogram-quality}
\end{figure}
\FloatBarrier
Quy trình thu tín hiệu được mô tả như sau:

\begin{figure}[htbp]
    \centering
    \begin{tikzpicture}[
        node distance=0.7cm and 0.55cm,
        box/.style={
            draw,
            rounded corners,
            minimum height=1cm,
            text width=2.45cm,
            align=center,
            fill=gray!8
        },
        arr/.style={-Latex, thick}
    ]

        \node[box] (rf) {Tín hiệu Drone/\\Controller};
        \node[box, right=of rf] (blade) {Bộ thu\\BladeRF};
        \node[box, right=of blade] (iq) {Tín hiệu số\\I/Q};
        \node[box, right=of iq] (pc) {Laptop\\đọc dữ liệu};
        \node[box, right=of pc] (store) {Lưu\\file I/Q};

        \draw[arr] (rf) -- (blade);
        \draw[arr] (blade) -- (iq);
        \draw[arr] (iq) -- (pc);
        \draw[arr] (pc) -- (store);

    \end{tikzpicture}
    \caption{Quy trình thu nhận tín hiệu RF sử dụng thiết bị BladeRF.}
    \label{fig:rf-capture}
\end{figure}


Trong mô hình này, anten của baldeRF thu nhận tín hiệu trong vùng phổ quan tâm, sau đó thiết bị SDR thực hiện khuếch đại, hạ tần về băng gốc, lọc theo băng thông cấu hình và lấy mẫu tương tự-số. Đầu ra của SDR là chuỗi mẫu phức:
\begin{equation}
    x[n] = I[n] + jQ[n],
\end{equation}
trong đó $I[n]$ và $Q[n]$ lần lượt là hai thành phần đồng pha và vuông pha của mẫu thứ $n$. Việc giữ cả hai thành phần này cho phép bảo toàn thông tin biên độ và pha của tín hiệu, phù hợp với phân tích phổ và tạo spectrogram.

Các tham số quan trọng của khối thu gồm tần số trung tâm $f_c$, tốc độ lấy mẫu $f_s$, băng thông thu, độ lợi thu, loại anten và định dạng lưu mẫu. Tần số trung tâm xác định vùng phổ được quan sát quanh băng tần drone. Tốc độ lấy mẫu và băng thông quyết định độ rộng phổ có thể ghi lại. Độ lợi thu ảnh hưởng trực tiếp đến mức công suất tín hiệu sau khuếch đại: nếu quá thấp, tín hiệu drone có thể bị chìm trong nhiễu nền; nếu quá cao, ADC (Analog-to-Digital Converter – bộ chuyển đổi tương tự sang số) có thể bị bão hòa và tạo artifact trên spectrogram.

Dữ liệu sau khi thu được lưu dưới dạng bản ghi I/Q liên tục. Mỗi bản ghi cần đi kèm thông tin cấu hình thu như $f_c$, $f_s$, băng thông, gain, thời điểm, loại tín hiệu và nhãn dự kiến. Metadata này không được đưa trực tiếp vào mô hình phân loại, nhưng cần thiết để đọc đúng file, tái tạo trục tần số, so sánh giữa các phiên thu và phân tích các trường hợp mô hình dự đoán sai.

Trong phạm vi đồ án, mô hình thu được sử dụng cho hai mục đích. Thứ nhất, hệ thống thu RF tạo dữ liệu \textit{Non-drone} để bổ sung lớp âm cho bài toán phân loại nhị phân, vì DroneDetectv2 không cung cấp đủ mẫu nền không chứa drone cho cấu hình huấn luyện của đồ án. Thứ hai, hệ thống thu được dùng để tạo các tập đánh giá độc lập nhằm kiểm tra khả năng tổng quát hóa khi điều kiện phần cứng và môi trường thu thay đổi.

\section{Tiền xử lý và phân đoạn dữ liệu}

Nguồn dữ liệu \textit{Drone} chính được sử dụng trong đồ án là bộ dữ liệu \textit{DroneDetectv2}. Do bộ dữ liệu này không cung cấp lớp \textit{Non-drone} tương ứng, đồ án tiến hành thu bổ sung các bản ghi RF không chứa hoạt động của drone bằng thiết bị bladeRF để xây dựng lớp \textit{Non-drone}. Sau khi thống nhất quy tắc gán nhãn, các bản ghi từ DroneDetectv2 được gán vào lớp \textit{Drone}, trong khi các bản ghi RF tự thu được gán vào lớp \textit{Non-drone}.

Mỗi bản ghi I/Q được quản lý như một đơn vị dữ liệu độc lập và được gán một mã định danh duy nhất. Sau quá trình tiền xử lý, toàn bộ các ảnh spectrogram sinh ra từ cùng một bản ghi được lưu trong một thư mục mang tên bản ghi tương ứng và kế thừa cùng một nhãn. Cấu trúc này giúp duy trì mối liên hệ giữa dữ liệu I/Q gốc và các ảnh spectrogram, đồng thời thuận tiện cho việc quản lý, truy vết và tái tạo dữ liệu khi cần thiết. Để hạn chế hiện tượng rò rỉ dữ liệu (\textit{data leakage}), việc chia tập huấn luyện, xác thực và kiểm thử được thực hiện ở mức bản ghi thay vì ở mức ảnh spectrogram. Theo đó, toàn bộ các ảnh được sinh ra từ một bản ghi chỉ xuất hiện trong duy nhất một tập dữ liệu, giúp kết quả đánh giá phản ánh đúng khả năng tổng quát hóa của mô hình. 

Mỗi bản ghi I/Q được chia thành các đoạn tín hiệu có độ dài cố định trước khi chuyển đổi thành ảnh spectrogram. Trong đồ án, mỗi đoạn có thời lượng $T=50$~ms và được xem là một mẫu đầu vào độc lập cho quá trình tạo dữ liệu. Nếu $f_s$ là tốc độ lấy mẫu, số mẫu của mỗi đoạn được xác định bởi:

\begin{equation}
L=T f_s.
\end{equation}

Đoạn thứ $k$ được biểu diễn bởi

\begin{equation}
\mathbf{x}_k=
\left[
x[kL],x[kL+1],\ldots,x[(k+1)L-1]
\right],
\end{equation}

trong đó các đoạn được lấy liên tiếp và không chồng lấn. Cách phân đoạn này giúp toàn bộ bản ghi được bao phủ bởi các khoảng thời gian liên tục, đồng thời đảm bảo mỗi ảnh spectrogram phản ánh trạng thái tín hiệu trong cùng một khoảng thời gian quan sát. Phần tín hiệu cuối bản ghi nếu không đủ $L$ mẫu sẽ được loại bỏ nhằm đảm bảo tất cả các đoạn có cùng độ dài. Mỗi đoạn tín hiệu sau đó được chuyển đổi thành một ảnh spectrogram với cùng bộ tham số STFT. Việc sử dụng cửa sổ thời gian cố định giúp các ảnh đầu vào có cùng kích thước, tạo điều kiện thuận lợi cho quá trình huấn luyện mô hình học sâu và đảm bảo tính nhất quán giữa các mẫu dữ liệu.

\section{Tạo ảnh spectrogram từ tín hiệu I/Q}

\subsection{Cấu hình STFT}

Dựa trên cơ sở STFT trình bày ở Chương~2, mỗi đoạn tín hiệu I/Q được biến đổi sang miền thời gian-tần số bằng một cấu hình xử lý thống nhất. Các tham số chính của STFT gồm loại cửa sổ (window), kích thước FFT ($N_{\mathrm{FFT}}$), khoảng dịch giữa hai cửa sổ liên tiếp (hop length) và phép dịch phổ (\texttt{fftshift}). Trong đồ án, cửa sổ Hamming được lựa chọn do có khả năng giảm hiện tượng rò rỉ phổ (spectral leakage) tốt hơn cửa sổ chữ nhật, đồng thời vẫn duy trì độ phân giải tần số phù hợp đối với tín hiệu RF. Kích thước FFT $N_{\mathrm{FFT}}$ quyết định độ phân giải theo miền tần số, với khoảng cách giữa hai bin tần số bằng $\Delta f=f_s/N_{\mathrm{FFT}}$, trong đó $f_s$ là tốc độ lấy mẫu. Giá trị $N_{\mathrm{FFT}}$ càng lớn thì độ phân giải tần số càng cao nhưng chi phí tính toán cũng tăng và độ phân giải thời gian giảm. Hop length xác định số mẫu dịch giữa hai cửa sổ STFT liên tiếp, qua đó quyết định mật độ các khung theo trục thời gian. Hop length nhỏ tạo ra nhiều khung hơn và giúp biểu diễn biến đổi theo thời gian mượt hơn, tuy nhiên làm tăng chi phí tính toán.

Sau khi tính FFT, phép \texttt{fftshift} được áp dụng để đưa thành phần tần số $0$ của tín hiệu băng gốc về giữa phổ. Khi đó, trục tần số của spectrogram được sắp xếp từ $-f_s/2$ đến gần $f_s/2$ quanh tần số trung tâm thu, giúp phổ đối xứng và thuận tiện cho việc quan sát các thành phần tín hiệu nằm ở cả hai phía của tần số trung tâm. Cách biểu diễn này cũng phù hợp với dữ liệu I/Q của SDR và thống nhất với phần cơ sở lý thuyết ở Chương~2. Các tham số STFT được cố định trong một cấu hình có thể tái lập. Với cùng một đoạn I/Q, thay đổi kích thước FFT, hop length hoặc quy tắc sắp xếp trục tần số có thể làm ảnh đầu vào khác đáng kể. Vì vậy, cấu hình STFT được xem như một phần của pipeline mô hình. Khi lưu checkpoint, cần lưu kèm cấu hình tạo spectrogram để bảo đảm giai đoạn suy luận tái tạo đúng dạng đầu vào mà mô hình đã học.

\subsection{Sinh và lưu ảnh spectrogram}

Sau khi phân đoạn, mỗi đoạn tín hiệu I/Q được chuyển đổi thành một ảnh spectrogram bằng phép biến đổi Fourier thời gian ngắn (STFT). Với đoạn tín hiệu $\mathbf{x}_k$, STFT tạo ra ma trận phổ phức $X[m,r]$, trong đó $m$ và $r$ lần lượt là chỉ số tần số và thời gian. Giá trị cường độ của spectrogram được tính theo thang logarit nhằm nén dải động của tín hiệu:

\begin{equation}
    S[m,r]
    =
    10\log_{10}
    \left(
        |X[m,r]|^2+\varepsilon
    \right),
\end{equation}

trong đó $\varepsilon$ là hằng số dương rất nhỏ nhằm tránh phép lấy logarit tại giá trị bằng không.

Ma trận $S$ biểu diễn sự phân bố năng lượng của tín hiệu theo thời gian và tần số. Trong đồ án, ma trận này được trực quan hóa bằng thư viện \texttt{Matplotlib} với bảng màu \texttt{jet}, trong đó các mức cường độ khác nhau được ánh xạ sang các màu khác nhau để tạo thành ảnh spectrogram. Sau khi sinh ảnh, tất cả các spectrogram được thay đổi kích thước về $224\times224$ pixel nhằm thống nhất kích thước đầu vào của mô hình học sâu.

Mỗi đoạn tín hiệu tương ứng với một ảnh spectrogram và được lưu thành một tệp riêng trong thư mục của bản ghi I/Q gốc. Việc lưu sẵn các ảnh spectrogram giúp tránh phải thực hiện lại phép biến đổi STFT trong mỗi lần huấn luyện hoặc đánh giá mô hình, từ đó rút ngắn thời gian tiền xử lý và đảm bảo toàn bộ tập dữ liệu được tạo theo cùng một quy trình.


\subsection{Chuẩn hóa đầu vào cho backbone}

Các ảnh spectrogram được lưu ở bước trước được đọc từ đĩa và sử dụng làm đầu vào cho backbone. Trong quá trình nạp dữ liệu, ảnh được chuyển về định dạng RGB và tiếp tục được áp dụng phép resize nhằm đảm bảo mọi ảnh đều có cùng kích thước đầu vào. Mặc dù spectrogram được tạo bằng bảng màu \texttt{jet}, ảnh sau đó được chuyển sang ảnh xám ba kênh. Cách xử lý này giúp giảm sự phụ thuộc của mô hình vào màu giả do colormap tạo ra, đồng thời vẫn đảm bảo tương thích với kiến trúc backbone tiền huấn luyện. Sau bước chuyển đổi kênh, ảnh được biểu diễn dưới dạng tensor và chuẩn hóa theo trung bình và độ lệch chuẩn của ImageNet:

\begin{equation}
    \widetilde{S}_{c}
    =
    \frac{S_{c}-\mu_c}{\sigma_c},
\end{equation}

trong đó $S_c$ là giá trị điểm ảnh của kênh thứ $c$, còn $\mu_c$ và $\sigma_c$ lần lượt là trung bình và độ lệch chuẩn sử dụng trong bước chuẩn hóa. Trong đồ án, các giá trị được sử dụng là
$\mu=[0.485,\,0.456,\,0.406]$ và
$\sigma=[0.229,\,0.224,\,0.225]$.

Toàn bộ pipeline tiền xử lý được áp dụng thống nhất cho tập huấn luyện, tập xác thực và tập kiểm thử. Đồ án không sử dụng các phép tăng cường dữ liệu hình học như xoay, lật hoặc dịch ảnh nhằm giữ nguyên đặc trưng vật lý của spectrogram và đảm bảo tính nhất quán giữa các giai đoạn huấn luyện và đánh giá mô hình.



\section{Xây dựng mô hình phân loại}

\subsection{Phân loại dựa trên trích xuất đặc trưng}

Phân loại dựa trên trích xuất đặc trưng hay \textit{Feature Extraction-based Classification} được dùng để đánh giá mức độ phù hợp của đặc trưng tiền huấn luyện khi backbone không được cập nhật. Với ảnh spectrogram $\mathbf{S}$, backbone đóng băng $g_{\phi}$ tạo embedding:
\begin{equation}
    \mathbf{h}=g_{\phi}(\mathbf{S}),
\end{equation}
trong đó $\phi$ là tham số tiền huấn luyện. Embedding có thể được lấy từ vector sau pooling, token đại diện hoặc lớp biểu diễn cuối tùy kiến trúc.

Các embedding của tập huấn luyện được dùng để huấn luyện SVM. Với SVM tuyến tính, hàm quyết định có dạng:

\begin{equation}
d(\mathbf{x})=\mathbf{w}^{T}\phi(\mathbf{x})+b,
\end{equation}

trong đó $\phi(\mathbf{x})$ là phép ánh xạ đặc trưng (đối với kernel SVM) và $b$ là hệ số dịch. Thay vì sử dụng ngưỡng mặc định $d(\mathbf{x})=0$, đồ án lựa chọn một ngưỡng phân loại $\tau$ trên tập xác thực nhằm tối ưu chỉ số đánh giá của mô hình. Quy tắc phân loại được xác định bởi

\begin{equation}
\hat{y}=
\begin{cases}
1, & d(\mathbf{x})\ge\tau,\\
0, & d(\mathbf{x})<\tau.
\end{cases}
\end{equation}

Ngưỡng $\tau$ sau khi xác định trên tập xác thực được giữ cố định và sử dụng cho toàn bộ quá trình đánh giá trên tập kiểm thử.

Ưu điểm của \textit{Feature Extraction-based Classification} là chi phí huấn luyện thấp và cho phép kiểm tra trực tiếp chất lượng embedding tiền huấn luyện. Hạn chế là backbone không thích nghi với miền spectrogram RF, nên các đặc trưng quan trọng như burst, dải phổ hoặc mẫu nhảy tần có thể chưa được biểu diễn tốt.

\subsection{Tinh chỉnh mô hình}
Tinh chỉnh mô hình hay \textit{Fine-tuning} cho phép cập nhật toàn bộ backbone theo bài toán phát hiện drone. Lớp phân loại gốc của mô hình được thay bằng lớp đầu ra cho hai lớp. Đồ án sử dụng hai logits với đầu ra của mô hình là:
\begin{equation}
    \mathbf{z}
    =
    [z_0,z_1]
    =
    f_{\theta}(\mathbf{S}),
\end{equation}
trong đó $z_0$ và $z_1$ lần lượt là điểm số của lớp \textit{Non-drone} và \textit{Drone}. Xác suất lớp được tính bằng softmax:
\begin{equation}
    p(y=c\mid\mathbf{S})
    =
    \frac{\exp(z_c)}
    {\exp(z_0)+\exp(z_1)},
    \qquad c\in\{0,1\}.
\end{equation}

Hàm mất mát cross-entropy được sử dụng trong quá trình huấn luyện:
\begin{equation}
    \mathcal{L}
    =
    -\frac{1}{B}
    \sum_{i=1}^{B}
    \sum_{c=0}^{1}
    \mathbb{I}(y_i=c)
    \log p(y=c\mid\mathbf{S}_i),
\end{equation}
trong đó $B$ là kích thước batch và $\mathbb{I}(\cdot)$ là hàm chỉ thị. Gradient được lan truyền qua lớp phân loại và backbone để cập nhật tham số.

\textit{Fine-tuning} giúp mô hình điều chỉnh đặc trưng theo cấu trúc riêng của spectrogram RF, nhưng có chi phí tính toán cao hơn và nguy cơ học cả đặc trưng riêng của nguồn dữ liệu huấn luyện. Vì vậy, mô hình được chọn bằng tập xác thực và được kiểm tra thêm trên dữ liệu RF tự thu.

\subsection{Backbone được khảo sát}

Đồ án khảo sát các backbone gồm VGG13-BN, ResNet50, EfficientNet-B2, ViT-B/16, Swin Transformer và DINOv2. Việc dùng nhiều backbone giúp đánh giá ảnh hưởng của kiến trúc và phương pháp tiền huấn luyện đối với miền ảnh spectrogram.

\begin{table}[ht]
\centering
\caption{Các backbone được khảo sát trong đồ án}
\label{tab:backbones}
\renewcommand{\arraystretch}{1.3}
\begin{tabular}{p{3cm} p{2.2cm} p{9.3cm}}
\hline
\textbf{Backbone} & \textbf{Kiến trúc} & \textbf{Đặc trưng nổi bật} \\
\hline

VGG13-BN &
CNN &
Kiến trúc CNN tuần tự gồm các lớp tích chập, Batch Normalization và pooling, có khả năng học hiệu quả các đặc trưng cục bộ của ảnh. \\

ResNet50 &
CNN &
Sử dụng các residual block và skip connection giúp giảm hiện tượng mất gradient, cho phép huấn luyện mạng sâu với hiệu quả cao. \\

EfficientNet-B2 &
CNN &
Áp dụng chiến lược compound scaling để cân bằng độ sâu, độ rộng và độ phân giải đầu vào, giúp nâng cao hiệu quả sử dụng tham số và chi phí tính toán. \\

ViT-B/16 &
Transformer &
Chia ảnh thành các patch kích thước $16\times16$ và sử dụng cơ chế self-attention để học mối quan hệ giữa các vùng ảnh trên phạm vi toàn cục. \\

Swin V2 Small &
Transformer &
Transformer phân cấp sử dụng window attention và shifted window để khai thác đồng thời đặc trưng cục bộ và toàn cục với chi phí tính toán thấp hơn Vision Transformer thông thường. \\

DINOv2 &
Transformer &
Backbone Vision Transformer được tiền huấn luyện bằng học tự giám sát, có khả năng học biểu diễn đặc trưng tổng quát và chuyển giao tốt sang các bài toán mới. \\

\hline
\end{tabular}
\end{table}

Tất cả backbone được đưa vào cùng quy trình tạo spectrogram, resize, chuẩn hóa và chia dữ liệu. Cách tổ chức này giúp kết quả phản ánh khác biệt của biểu diễn mô hình và chiến lược huấn luyện, thay vì khác biệt do tiền xử lý.

\subsection{Quy trình huấn luyện và suy luận}

\subsection*{Huấn luyện và lựa chọn mô hình}

Tập dữ liệu được chia thành tập huấn luyện, tập xác thực và tập kiểm thử. Tập huấn luyện dùng để tối ưu SVM trong \textit{Feature Extraction-based Classification} hoặc cập nhật tham số trong \textit{Fine-tuning}. Tập xác thực dùng để chọn siêu tham số, checkpoint và ngưỡng quyết định. Tập kiểm thử chỉ được sử dụng sau khi quá trình lựa chọn mô hình hoàn tất.

Trong \textit{Feature Extraction-based Classification}, spectrogram được đưa qua backbone đóng băng để trích xuất embedding. Các embedding được chuẩn hóa theo thống kê của tập huấn luyện rồi đưa vào SVM. Trong \textit{fine-tuning}, ảnh spectrogram được đưa trực tiếp qua mô hình, lớp phân loại sinh logits và mô hình được tối ưu bằng cross-entropy. Các cấu hình huấn luyện chi tiết như learning rate, batch size, số epoch và tiêu chí chọn checkpoint được trình bày ở chương thực nghiệm.

Checkpoint được lựa chọn dựa trên kết quả của tập xác thực thay vì tập kiểm thử. Cách làm này giúp tập kiểm thử giữ vai trò đánh giá độc lập. Trong quá trình huấn luyện, các thông tin như loss, accuracy, learning rate và checkpoint tốt nhất cần được lưu lại để phục vụ tái lập thí nghiệm và so sánh giữa các backbone.

\subsection*{Suy luận trên tập kiểm thử nội miền}

Suy luận nội miền được thực hiện trên tập kiểm thử nhị phân không tham gia huấn luyện hoặc lựa chọn mô hình, gồm các mẫu \textit{Drone} từ DroneDetectv2 và các mẫu \textit{Non-drone} tự thu, được xử lý bằng cùng pipeline như dữ liệu huấn luyện. Mỗi bản ghi được phân đoạn, chuyển thành spectrogram và đưa vào mô hình. Với \textit{Feature Extraction-based Classification}, embedding được trích xuất và phân loại bằng SVM; với \textit{Fine-tuning}, mô hình dự đoán trực tiếp nhãn có xác suất lớn nhất. Kết quả suy luận được sử dụng để tính các chỉ số đánh giá như accuracy, precision, recall, F1-score và ma trận nhầm lẫn.

\subsection*{Suy luận trên dữ liệu RF tự thu}

Đối với dữ liệu RF tự thu, mỗi bản ghi I/Q được phân đoạn và chuyển thành spectrogram bằng cùng cấu hình tiền xử lý như dữ liệu huấn luyện. Các tập dữ liệu này chỉ được sử dụng để đánh giá, không tham gia huấn luyện hoặc lựa chọn mô hình. Do dữ liệu tự thu có thể khác dữ liệu nội miền về thiết bị SDR, anten, mức gain, môi trường thu và điều kiện tín hiệu, kết quả suy luận phản ánh khả năng tổng quát hóa của mô hình khi xuất hiện hiện tượng \textit{domain shift}, đồng thời đánh giá tính khả thi của hệ thống trong điều kiện triển khai thực tế.


Chương này đã trình bày phương pháp đề xuất cho bài toán phát hiện drone bằng tín hiệu RF. Pipeline bắt đầu từ mô hình thu tín hiệu RF bằng SDR, lưu dữ liệu I/Q, phân đoạn, tạo spectrogram, chuẩn hóa ảnh và đưa vào mô hình phân loại hai lớp \textit{Drone} / \textit{Non-drone}. Các bước xử lý được thiết kế thống nhất để hạn chế sai lệch giữa dữ liệu huấn luyện, dữ liệu kiểm thử và dữ liệu RF tự thu. Hai chiến lược khai thác mô hình tiền huấn luyện được xây dựng gồm \textit{Feature Extraction-based Classification} với SVM trên embedding đóng băng và \textit{fine-tuning} toàn bộ backbone bằng cross-entropy. Sáu backbone gồm VGG13-BN, ResNet50, EfficientNet-B2, ConvNeXtV2, Swin Transformer và DINOv2 được đặt trong cùng pipeline để so sánh công bằng hơn. Quy trình này tạo cơ sở để chương tiếp theo đánh giá hiệu năng nội miền, khả năng tổng quát hóa khi môi trường thay đổi

\end{document}
